\documentclass[11pt]{article}
\usepackage[margin=1in]{geometry}
\usepackage{amsmath, amssymb}
\usepackage{graphicx}
\usepackage{setspace}
\usepackage{hyperref}

\title{One-Class SVM for PFAS--MOF Descriptor Modeling}
\author{}
\date{}

\begin{document}
\maketitle

\section{Problem framing}

\begin{itemize}
\item \textbf{Goal:} Learn the \emph{region} of descriptor space occupied by known PFAS-adsorbing MOFs using positive-only data.
\item \textbf{Why OC-SVM:} No reliable negatives or scalar target are available; instead, a soft boundary around positive examples is required.
\end{itemize}

\section{Descriptor set}

\subsection*{Core (chemist-prioritized)}
\begin{itemize}
\item PLD (\AA)
\item LCD (\AA)
\item Geometric Surface Area (m$^2$/g)
\item Void fraction
\item Pore volume (cm$^3$/g)
\end{itemize}

\subsection*{Candidate / exploratory}
\begin{itemize}
\item Density (g/cm$^3$)
\end{itemize}

\noindent\textbf{Dimensionality:} A feature space of 5--6 descriptors is appropriate for OC-SVM given sample sizes of 23 (PFOA) and 13 (PFOS).

\section{Modeling choices}

\subsection{What the OC-SVM is optimizing (intuition)}

\begin{itemize}
\item The one-class SVM learns a \emph{compact region} of descriptor space that contains most of the positive data.
\item Unlike binary SVMs, there is no opposing class pushing the boundary from the outside.
\item The solution is determined by a trade-off between:
\begin{itemize}
\item minimizing the volume (or norm) of the enclosing region, and
\item allowing a controlled fraction of points to lie outside the boundary.
\end{itemize}
\item This trade-off is governed primarily by the parameter $\nu$, which prevents the trivial solution of an infinitely large enclosing region.
\item A single multivariate OC-SVM is standard; using more than two features is routine (two-dimensional examples are used only for visualization).
\end{itemize}

\subsection{Kernel choice (justification)}

The \textbf{Radial Basis Function (RBF)} kernel is preferred for one-class problems in materials descriptor spaces.

\paragraph{Geometric rationale.}
\begin{itemize}
\item The OC-SVM objective is to \emph{enclose} the region occupied by positive examples.
\item RBF kernels induce closed, flexible boundaries in feature space that can wrap around clusters of points.
\item This matches the scientific goal of learning a compact support of PFAS-compatible MOFs.
\end{itemize}

\noindent\textbf{Kernel:} RBF.
\begin{itemize}
\item $\nu$: controls the allowed fraction of outliers and acts as a regularization parameter preventing unbounded expansion of the region.
\item $\gamma$: sets the characteristic length scale of the boundary and is chosen via simple heuristics (e.g., median pairwise distance).
\end{itemize}

\section{Critical preprocessing}

\begin{itemize}
\item \textbf{Scaling is mandatory.} Z-score or robust scaling must be applied.
\item The RBF kernel relies on Euclidean distances; unscaled features (e.g., surface area) would dominate the geometry.
\item Correlated descriptors (e.g., PLD, LCD, pore volume) are acceptable, as the OC-SVM operates on the induced distance geometry rather than assuming feature independence.
\item PCA is used strictly for visualization and interpretation and is not required for model training.
\end{itemize}

\section{Assessing descriptor influence and sensitivity}

There is no native coefficient-based notion of feature importance. Instead, descriptor relevance is assessed using \textbf{sensitivity analyses}.

\subsection{Leave-One-Feature-Out (LOFO)}
\begin{itemize}
\item Retrain the model after removing each descriptor in turn.
\item Compare the fraction of points flagged as outliers, changes in anomaly score distributions, and boundary stability.
\end{itemize}

\subsection{Perturbation sensitivity}
\begin{itemize}
\item Perturb a single descriptor (e.g., $\pm 5$--$10\%$) while holding others fixed.
\item Observe changes in anomaly scores; this is often interpretable for domain experts.
\end{itemize}

\section{What not to do}

\begin{itemize}
\item Do not report p-values or regression-style coefficients.
\item Do not assume monotonic or linear effects.
\item Do not over-tune hyperparameters to force perceived importance.
\end{itemize}

\section{Visualizing the learned region}

\noindent\emph{Note:} Different choices of slice (e.g., PCA planes versus original feature-pair slices, or different fixed reference values) may lead to different apparent boundaries. Such variation reflects the high-dimensional nature of the learned region rather than model instability.

\begin{itemize}
\item Any two-dimensional decision boundary plot represents a \textbf{conditional slice} of the learned high-dimensional region.
\item PCA-based plots show the intersection of the decision function with the PCA subspace.
\item Original-feature plots vary two selected descriptors while fixing all remaining features at their mean values.
\item These visualizations are illustrative and should not be interpreted as complete projections of the full decision boundary.
\end{itemize}

\section{Recommended workflow}

\begin{enumerate}
\item Train an RBF OC-SVM on the scaled core descriptors.
\item Validate stability with respect to $\nu$.
\item Add density and compare resulting boundaries and scores.
\item Perform LOFO and perturbation analyses.
\item Report boundary stability, sensitivity, and interpretable contour-based visualizations (PCA or feature-pair slices), not coefficients.
\end{enumerate}

\section{Outputs to report}

\begin{itemize}
\item Continuous anomaly scores.
\item In/out decisions at the chosen threshold.
\item Sensitivity results supporting descriptor relevance.
\item Visualization figures.
\end{itemize}

\end{document}